% Options for packages loaded elsewhere
\PassOptionsToPackage{unicode}{hyperref}
\PassOptionsToPackage{hyphens}{url}
\documentclass[
  11pt,
  a4paper,
]{article}
\usepackage{xcolor}
\usepackage[margin=18mm]{geometry}
\usepackage{amsmath,amssymb}
\setcounter{secnumdepth}{-\maxdimen} % remove section numbering
\usepackage{iftex}
\ifPDFTeX
  \usepackage[T1]{fontenc}
  \usepackage[utf8]{inputenc}
  \usepackage{textcomp} % provide euro and other symbols
\else % if luatex or xetex
  \usepackage{unicode-math} % this also loads fontspec
  \defaultfontfeatures{Scale=MatchLowercase}
  \defaultfontfeatures[\rmfamily]{Ligatures=TeX,Scale=1}
\fi
\ifPDFTeX\else
  % xetex/luatex font selection
    \setmainfont[]{Noto Serif}
    \setsansfont[]{Noto Sans}
    \setmonofont[]{Noto Sans Mono}
  \setmathfont[]{Noto Sans Math}
\fi
% Use upquote if available, for straight quotes in verbatim environments
\IfFileExists{upquote.sty}{\usepackage{upquote}}{}
\IfFileExists{microtype.sty}{% use microtype if available
  \usepackage[]{microtype}
  \UseMicrotypeSet[protrusion]{basicmath} % disable protrusion for tt fonts
}{}
\makeatletter
\@ifundefined{KOMAClassName}{% if non-KOMA class
  \IfFileExists{parskip.sty}{%
    \usepackage{parskip}
  }{% else
    \setlength{\parindent}{0pt}
    \setlength{\parskip}{6pt plus 2pt minus 1pt}}
}{% if KOMA class
  \KOMAoptions{parskip=half}}
\makeatother
\ifLuaTeX
\usepackage[bidi=basic]{babel}
\else
\usepackage[bidi=default]{babel}
\fi
\babelprovide[main,import]{russian}
\ifPDFTeX
\else
\babelfont{rm}[]{Noto Serif}
\fi
% get rid of language-specific shorthands (see #6817):
\let\LanguageShortHands\languageshorthands
\def\languageshorthands#1{}
\setlength{\emergencystretch}{3em} % prevent overfull lines
\providecommand{\tightlist}{%
  \setlength{\itemsep}{0pt}\setlength{\parskip}{0pt}}
\usepackage{bookmark}
\IfFileExists{xurl.sty}{\usepackage{xurl}}{} % add URL line breaks if available
\urlstyle{same}
\hypersetup{
  pdftitle={Билеты по дискретным случайным процессам},
  pdflang={ru-RU},
  hidelinks,
  pdfcreator={LaTeX via pandoc}}

\title{Билеты по дискретным случайным процессам}
\author{}
\date{}

\begin{document}
\maketitle

\section{Билет 09. Теоремы Фубини и независимые случайные
величины}\label{ux431ux438ux43bux435ux442-09.-ux442ux435ux43eux440ux435ux43cux44b-ux444ux443ux431ux438ux43dux438-ux438-ux43dux435ux437ux430ux432ux438ux441ux438ux43cux44bux435-ux441ux43bux443ux447ux430ux439ux43dux44bux435-ux432ux435ux43bux438ux447ux438ux43dux44b}

Источники: Ширяев 1, гл. I, §3; гл. II, §§5-6; лекции ДСП.

\subsection{1. Короткий
ответ}\label{ux43aux43eux440ux43eux442ux43aux438ux439-ux43eux442ux432ux435ux442}

В этом билете две связанные идеи.

Первая: если на произведении измеримых пространств задана произведенная
мера, то интеграл по произведению можно считать как повторный интеграл.
Для неотрицательной измеримой функции работает теорема Тонелли:

\[
\int_{\Omega_1\times\Omega_2} f(\omega_1,\omega_2)\,(\mu_1\otimes\mu_2)(d\omega_1,d\omega_2)
=
\int_{\Omega_1}
\left(
\int_{\Omega_2} f(\omega_1,\omega_2)\,\mu_2(d\omega_2)
\right)
\mu_1(d\omega_1),
\]

и аналогично с другим порядком интегрирования. Значения могут быть
бесконечными.

Вторая: независимость означает, что совместное распределение
раскладывается в произведение маргинальных распределений. Для случайных
элементов \(X\) и \(Y\):

\[
X\ \text{и}\ Y\ \text{независимы}
\quad\Longleftrightarrow\quad
P_{(X,Y)}=P_X\otimes P_Y.
\]

Отсюда через Фубини получается ключевая формула:

\[
\mathbb E[g(X)h(Y)]
=
\mathbb E g(X)\,\mathbb E h(Y),
\]

если соответствующие интегралы корректно определены. Для независимых
случайных величин с плотностями совместная плотность равна произведению:

\[
p_{X,Y}(x,y)=p_X(x)p_Y(y)
\]

почти всюду.

\subsection{2. Полный
ответ}\label{ux43fux43eux43bux43dux44bux439-ux43eux442ux432ux435ux442}

\subsubsection{2.1. Введение и область
применения}\label{ux432ux432ux435ux434ux435ux43dux438ux435-ux438-ux43eux431ux43bux430ux441ux442ux44c-ux43fux440ux438ux43cux435ux43dux435ux43dux438ux44f}

Теоремы Фубини нужны для того, чтобы корректно переходить от интеграла
по произведению пространств к повторным интегралам. В вероятности это
появляется постоянно: совместные распределения случайных векторов,
математические ожидания функций нескольких случайных величин, свертки
распределений, марковские переходы, корреляционные функции и
спектральные формулы.

\subsubsection{2.2. Конструкция произведения пространств и
мер}\label{ux43aux43eux43dux441ux442ux440ux443ux43aux446ux438ux44f-ux43fux440ux43eux438ux437ux432ux435ux434ux435ux43dux438ux44f-ux43fux440ux43eux441ux442ux440ux430ux43dux441ux442ux432-ux438-ux43cux435ux440}

Пусть заданы два измеримых пространства с мерами:

\[
(\Omega_1,\mathcal F_1,\mu_1),
\qquad
(\Omega_2,\mathcal F_2,\mu_2).
\]

На декартовом произведении \(\Omega_1\times\Omega_2\) рассматривают
произведенную \(\sigma\)-алгебру

\[
\mathcal F_1\otimes\mathcal F_2
=
\sigma\{A\times B:\ A\in\mathcal F_1,\ B\in\mathcal F_2\}.
\]

Если меры \(\mu_1\) и \(\mu_2\) \(\sigma\)-конечны, то существует
единственная мера \(\mu_1\otimes\mu_2\) на
\(\mathcal F_1\otimes\mathcal F_2\), такая что

\[
(\mu_1\otimes\mu_2)(A\times B)=\mu_1(A)\mu_2(B).
\]

Для вероятностных мер \(\sigma\)-конечность автоматически выполнена,
потому что \(P(\Omega)=1\).

\subsubsection{2.3. Сечения функций и теорема
Тонелли}\label{ux441ux435ux447ux435ux43dux438ux44f-ux444ux443ux43dux43aux446ux438ux439-ux438-ux442ux435ux43eux440ux435ux43cux430-ux442ux43eux43dux435ux43bux43bux438}

Пусть

\[
f:\Omega_1\times\Omega_2\to[0,\infty]
\]

измерима относительно \(\mathcal F_1\otimes\mathcal F_2\). Тогда для
фиксированного \(\omega_1\) функция

\[
\omega_2\mapsto f(\omega_1,\omega_2)
\]

измерима по \(\mathcal F_2\), а для фиксированного \(\omega_2\) функция

\[
\omega_1\mapsto f(\omega_1,\omega_2)
\]

измерима по \(\mathcal F_1\). Эти функции называются сечениями функции
\(f\).

Теорема Тонелли утверждает, что для неотрицательной измеримой функции
повторные интегралы существуют как расширенные неотрицательные числа и
равны интегралу по произведенной мере:

\[
\int f\,d(\mu_1\otimes\mu_2)
=
\int_{\Omega_1}
\left(\int_{\Omega_2} f(\omega_1,\omega_2)\,\mu_2(d\omega_2)\right)
\mu_1(d\omega_1)
\]

\[
=
\int_{\Omega_2}
\left(\int_{\Omega_1} f(\omega_1,\omega_2)\,\mu_1(d\omega_1)\right)
\mu_2(d\omega_2).
\]

Здесь не требуется заранее знать конечность интеграла. Можно получить
значение \(+\infty\). Это важно: для неотрицательных функций порядок
интегрирования можно менять без проверки абсолютной интегрируемости.

\subsubsection{2.4. Теорема Фубини для интегрируемых
функций}\label{ux442ux435ux43eux440ux435ux43cux430-ux444ux443ux431ux438ux43dux438-ux434ux43bux44f-ux438ux43dux442ux435ux433ux440ux438ux440ux443ux435ux43cux44bux445-ux444ux443ux43dux43aux446ux438ux439}

Теорема Фубини - это версия для знакопеременных или комплекснозначных
функций. Пусть

\[
f\in L^1(\mu_1\otimes\mu_2),
\qquad
\int_{\Omega_1\times\Omega_2}|f|\,d(\mu_1\otimes\mu_2)<\infty.
\]

Тогда:

\begin{enumerate}
\def\labelenumi{\arabic{enumi}.}
\tightlist
\item
  для \(\mu_1\)-почти всех \(\omega_1\) функция \(f(\omega_1,\cdot)\)
  интегрируема по \(\mu_2\);
\item
  для \(\mu_2\)-почти всех \(\omega_2\) функция \(f(\cdot,\omega_2)\)
  интегрируема по \(\mu_1\);
\item
  функции повторных интегралов измеримы и интегрируемы;
\item
  выполнено равенство
\end{enumerate}

\[
\int_{\Omega_1\times\Omega_2} f\,d(\mu_1\otimes\mu_2)
=
\int_{\Omega_1}
\left(\int_{\Omega_2} f(\omega_1,\omega_2)\,\mu_2(d\omega_2)\right)
\mu_1(d\omega_1)
\]

\[
=
\int_{\Omega_2}
\left(\int_{\Omega_1} f(\omega_1,\omega_2)\,\mu_1(d\omega_1)\right)
\mu_2(d\omega_2).
\]

Главное условие Фубини - интегрируемость \(|f|\) по произведенной мере.
Практически часто проверяют достаточное условие:

\[
\int_{\Omega_1}
\left(\int_{\Omega_2}|f(\omega_1,\omega_2)|\,\mu_2(d\omega_2)\right)
\mu_1(d\omega_1)
<\infty.
\]

Тогда по Тонелли

\[
\int |f|\,d(\mu_1\otimes\mu_2)<\infty,
\]

значит применима теорема Фубини.

\subsubsection{2.5. Независимость событий и
сигма-алгебр}\label{ux43dux435ux437ux430ux432ux438ux441ux438ux43cux43eux441ux442ux44c-ux441ux43eux431ux44bux442ux438ux439-ux438-ux441ux438ux433ux43cux430-ux430ux43bux433ux435ux431ux440}

Теперь перейдем к независимости. Для двух событий \(A,B\in\mathcal F\)
независимость определяется формулой

\[
P(A\cap B)=P(A)P(B).
\]

Интуиция: знание того, что произошло \(A\), не меняет вероятность \(B\).
Если \(P(A)>0\), это равносильно

\[
P(B\mid A)=P(B).
\]

Для нескольких событий \(A_1,\ldots,A_n\) независимость в совокупности
означает, что для любого непустого набора индексов \(i_1,\ldots,i_k\)

\[
P(A_{i_1}\cap\cdots\cap A_{i_k})
=
P(A_{i_1})\cdots P(A_{i_k}).
\]

Важно: попарная независимость не равна независимости в совокупности.

Для \(\sigma\)-алгебр
\(\mathcal G_1,\ldots,\mathcal G_n\subset\mathcal F\) независимость
означает:

\[
P(G_1\cap\cdots\cap G_n)
=
P(G_1)\cdots P(G_n)
\]

для любых \(G_i\in\mathcal G_i\). Случайные элементы \(X_1,\ldots,X_n\)
независимы, если независимы порожденные ими \(\sigma\)-алгебры:

\[
\sigma(X_1),\ldots,\sigma(X_n).
\]

Эквивалентно, для любых борелевских множеств \(B_i\):

\[
P\{X_1\in B_1,\ldots,X_n\in B_n\}
=
\prod_{i=1}^n P\{X_i\in B_i\}.
\]

\subsubsection{2.6. Независимость случайных элементов и совместное
распределение}\label{ux43dux435ux437ux430ux432ux438ux441ux438ux43cux43eux441ux442ux44c-ux441ux43bux443ux447ux430ux439ux43dux44bux445-ux44dux43bux435ux43cux435ux43dux442ux43eux432-ux438-ux441ux43eux432ux43cux435ux441ux442ux43dux43eux435-ux440ux430ux441ux43fux440ux435ux434ux435ux43bux435ux43dux438ux435}

В терминах распределений это записывается особенно компактно. Для двух
случайных элементов

\[
X:\Omega\to E,
\qquad
Y:\Omega\to F
\]

с распределениями \(P_X\) и \(P_Y\) совместное распределение вектора
\((X,Y)\) есть мера

\[
P_{(X,Y)}(C)=P\{(X,Y)\in C\},
\qquad
C\in\mathcal E\otimes\mathcal F.
\]

Тогда

\[
X\ \text{и}\ Y\ \text{независимы}
\quad\Longleftrightarrow\quad
P_{(X,Y)}=P_X\otimes P_Y.
\]

Действительно, на прямоугольниках \(A\times B\) равенство дает

\[
P\{X\in A,\ Y\in B\}
=
P_X(A)P_Y(B),
\]

то есть ровно определение независимости. А поскольку прямоугольники
порождают произведенную \(\sigma\)-алгебру, равенство переносится на все
измеримые множества стандартным аргументом через монотонный класс.

\subsubsection{2.7. Функции распределения и
плотности}\label{ux444ux443ux43dux43aux446ux438ux438-ux440ux430ux441ux43fux440ux435ux434ux435ux43bux435ux43dux438ux44f-ux438-ux43fux43bux43eux442ux43dux43eux441ux442ux438}

Если \(X\) и \(Y\) - вещественные случайные величины, то независимость
эквивалентна факторизации совместной функции распределения:

\[
F_{X,Y}(x,y)
=
P\{X\le x,\ Y\le y\}
=
F_X(x)F_Y(y)
\]

для всех \(x,y\in\mathbb R\).

Если у совместного распределения есть плотность \(p_{X,Y}\) относительно
меры Лебега на \(\mathbb R^2\), то маргинальные плотности выражаются
через повторные интегралы:

\[
p_X(x)=\int_{\mathbb R}p_{X,Y}(x,y)\,dy,
\qquad
p_Y(y)=\int_{\mathbb R}p_{X,Y}(x,y)\,dx.
\]

Здесь используется Фубини/Тонелли. В этом случае

\[
X\ \text{и}\ Y\ \text{независимы}
\quad\Longleftrightarrow\quad
p_{X,Y}(x,y)=p_X(x)p_Y(y)
\]

почти всюду относительно двумерной меры Лебега.

\subsubsection{2.8. Математическое ожидание произведения
функций}\label{ux43cux430ux442ux435ux43cux430ux442ux438ux447ux435ux441ux43aux43eux435-ux43eux436ux438ux434ux430ux43dux438ux435-ux43fux440ux43eux438ux437ux432ux435ux434ux435ux43dux438ux44f-ux444ux443ux43dux43aux446ux438ux439}

Фубини дает основную вычислительную формулу для независимых величин.
Пусть \(X\) и \(Y\) независимы, а \(g\) и \(h\) - измеримые функции.
Если \(g(X)\) и \(h(Y)\) неотрицательны или если все нужные величины
интегрируемы, то

\[
\mathbb E[g(X)h(Y)]
=
\int_{E\times F} g(x)h(y)\,P_{(X,Y)}(dx,dy).
\]

По независимости

\[
P_{(X,Y)}=P_X\otimes P_Y.
\]

По Фубини:

\[
\mathbb E[g(X)h(Y)]
=
\int_E\int_F g(x)h(y)\,P_Y(dy)\,P_X(dx)
\]

\[
=
\left(\int_E g(x)\,P_X(dx)\right)
\left(\int_F h(y)\,P_Y(dy)\right)
=
\mathbb E g(X)\,\mathbb E h(Y).
\]

В частности, если \(X,Y\in L^1\) и независимы, то

\[
\mathbb E[XY]=\mathbb EX\,\mathbb EY.
\]

Если дополнительно \(X,Y\in L^2\), то

\[
\operatorname{Cov}(X,Y)=0.
\]

Но обратное неверно: нулевая ковариация вообще не означает
независимости. Исключение - специальные классы, например совместно
гауссовские величины, что относится к Билет 18.

\subsubsection{2.9. Распределение суммы независимых величин
(свертка)}\label{ux440ux430ux441ux43fux440ux435ux434ux435ux43bux435ux43dux438ux435-ux441ux443ux43cux43cux44b-ux43dux435ux437ux430ux432ux438ux441ux438ux43cux44bux445-ux432ux435ux43bux438ux447ux438ux43d-ux441ux432ux435ux440ux442ux43aux430}

Еще одно важное следствие - распределение суммы независимых величин.
Если \(X\) и \(Y\) независимы, то закон \(X+Y\) равен свертке законов:

\[
P_{X+Y}=P_X*P_Y.
\]

Для борелевского множества \(B\subset\mathbb R\):

\[
P\{X+Y\in B\}
=
\int_{\mathbb R}\int_{\mathbb R}\mathbf 1_B(x+y)\,P_X(dx)P_Y(dy).
\]

В дискретном случае:

\[
P\{X+Y=z\}
=
\sum_x P\{X=x\}P\{Y=z-x\}.
\]

Если есть плотности \(p_X\) и \(p_Y\), то плотность суммы:

\[
p_{X+Y}(z)
=
\int_{\mathbb R}p_X(x)p_Y(z-x)\,dx.
\]

Через характеристические функции это превращается в произведение:

\[
\varphi_{X+Y}(t)=\varphi_X(t)\varphi_Y(t),
\]

что связывает билет с Билет 08 и предельными теоремами.

\subsubsection{2.10. Резюме}\label{ux440ux435ux437ux44eux43cux435}

Итак, логика билета такая: произведенная мера описывает независимое
совместное распределение; Фубини и Тонелли позволяют корректно
интегрировать по этому произведению; поэтому из независимости следуют
факторизация ожиданий, произведение плотностей, свертка распределений и
произведение характеристических функций.

\subsection{3. Основные
определения}\label{ux43eux441ux43dux43eux432ux43dux44bux435-ux43eux43fux440ux435ux434ux435ux43bux435ux43dux438ux44f}

\subsubsection{\texorpdfstring{Произведенная
\(\sigma\)-алгебра}{Произведенная \textbackslash sigma-алгебра}}\label{ux43fux440ux43eux438ux437ux432ux435ux434ux435ux43dux43dux430ux44f-sigma-ux430ux43bux433ux435ux431ux440ux430}

Для измеримых пространств \((\Omega_1,\mathcal F_1)\) и
\((\Omega_2,\mathcal F_2)\):

\[
\mathcal F_1\otimes\mathcal F_2
=
\sigma\{A\times B:\ A\in\mathcal F_1,\ B\in\mathcal F_2\}.
\]

\subsubsection{Произведение
мер}\label{ux43fux440ux43eux438ux437ux432ux435ux434ux435ux43dux438ux435-ux43cux435ux440}

Если \(\mu_1\) и \(\mu_2\) \(\sigma\)-конечны, то \(\mu_1\otimes\mu_2\)
- мера на \(\mathcal F_1\otimes\mathcal F_2\), задаваемая на
прямоугольниках формулой

\[
(\mu_1\otimes\mu_2)(A\times B)=\mu_1(A)\mu_2(B).
\]

\subsubsection{Сечение
множества}\label{ux441ux435ux447ux435ux43dux438ux435-ux43cux43dux43eux436ux435ux441ux442ux432ux430}

Для \(C\subset\Omega_1\times\Omega_2\):

\[
C_{\omega_1}=\{\omega_2:\ (\omega_1,\omega_2)\in C\},
\qquad
C^{\omega_2}=\{\omega_1:\ (\omega_1,\omega_2)\in C\}.
\]

\subsubsection{Сечение
функции}\label{ux441ux435ux447ux435ux43dux438ux435-ux444ux443ux43dux43aux446ux438ux438}

Для \(f(\omega_1,\omega_2)\):

\[
f_{\omega_1}(\omega_2)=f(\omega_1,\omega_2),
\qquad
f^{\omega_2}(\omega_1)=f(\omega_1,\omega_2).
\]

\subsubsection{Интегрируемость на
произведении}\label{ux438ux43dux442ux435ux433ux440ux438ux440ux443ux435ux43cux43eux441ux442ux44c-ux43dux430-ux43fux440ux43eux438ux437ux432ux435ux434ux435ux43dux438ux438}

Функция \(f\) интегрируема относительно \(\mu_1\otimes\mu_2\), если

\[
\int_{\Omega_1\times\Omega_2}|f|\,d(\mu_1\otimes\mu_2)<\infty.
\]

\subsubsection{Независимые
события}\label{ux43dux435ux437ux430ux432ux438ux441ux438ux43cux44bux435-ux441ux43eux431ux44bux442ux438ux44f}

События \(A\) и \(B\) независимы, если

\[
P(A\cap B)=P(A)P(B).
\]

\subsubsection{Независимость событий в
совокупности}\label{ux43dux435ux437ux430ux432ux438ux441ux438ux43cux43eux441ux442ux44c-ux441ux43eux431ux44bux442ux438ux439-ux432-ux441ux43eux432ux43eux43aux443ux43fux43dux43eux441ux442ux438}

События \(A_1,\ldots,A_n\) независимы в совокупности, если для любого
непустого набора индексов \(i_1,\ldots,i_k\):

\[
P(A_{i_1}\cap\cdots\cap A_{i_k})
=
P(A_{i_1})\cdots P(A_{i_k}).
\]

\subsubsection{\texorpdfstring{Независимые
\(\sigma\)-алгебры}{Независимые \textbackslash sigma-алгебры}}\label{ux43dux435ux437ux430ux432ux438ux441ux438ux43cux44bux435-sigma-ux430ux43bux433ux435ux431ux440ux44b}

\(\sigma\)-алгебры \(\mathcal G_1,\ldots,\mathcal G_n\) независимы, если
для любых \(G_i\in\mathcal G_i\):

\[
P(G_1\cap\cdots\cap G_n)
=
\prod_{i=1}^nP(G_i).
\]

\subsubsection{Независимые случайные
элементы}\label{ux43dux435ux437ux430ux432ux438ux441ux438ux43cux44bux435-ux441ux43bux443ux447ux430ux439ux43dux44bux435-ux44dux43bux435ux43cux435ux43dux442ux44b}

Случайные элементы \(X_1,\ldots,X_n\) независимы, если независимы
\(\sigma(X_1),\ldots,\sigma(X_n)\).

\subsubsection{Распределительный критерий
независимости}\label{ux440ux430ux441ux43fux440ux435ux434ux435ux43bux438ux442ux435ux43bux44cux43dux44bux439-ux43aux440ux438ux442ux435ux440ux438ux439-ux43dux435ux437ux430ux432ux438ux441ux438ux43cux43eux441ux442ux438}

\[
X_1,\ldots,X_n\ \text{независимы}
\quad\Longleftrightarrow\quad
P_{(X_1,\ldots,X_n)}
=
P_{X_1}\otimes\cdots\otimes P_{X_n}.
\]

\subsubsection{Свертка
распределений}\label{ux441ux432ux435ux440ux442ux43aux430-ux440ux430ux441ux43fux440ux435ux434ux435ux43bux435ux43dux438ux439}

Для вероятностных мер \(\mu\) и \(\nu\) на \(\mathbb R\) свертка
\(\mu*\nu\) задается формулой

\[
(\mu*\nu)(B)
=
\int_{\mathbb R}\int_{\mathbb R}\mathbf 1_B(x+y)\,\mu(dx)\nu(dy).
\]

Если \(X,Y\) независимы, то

\[
P_{X+Y}=P_X*P_Y.
\]

\subsection{4. Основные теоремы и
свойства}\label{ux43eux441ux43dux43eux432ux43dux44bux435-ux442ux435ux43eux440ux435ux43cux44b-ux438-ux441ux432ux43eux439ux441ux442ux432ux430}

\subsubsection{Теорема 1. Существование произведения
мер}\label{ux442ux435ux43eux440ux435ux43cux430-1.-ux441ux443ux449ux435ux441ux442ux432ux43eux432ux430ux43dux438ux435-ux43fux440ux43eux438ux437ux432ux435ux434ux435ux43dux438ux44f-ux43cux435ux440}

Если \(\mu_1\) и \(\mu_2\) - \(\sigma\)-конечные меры, то существует
единственная мера \(\mu_1\otimes\mu_2\) на
\(\mathcal F_1\otimes\mathcal F_2\), такая что

\[
(\mu_1\otimes\mu_2)(A\times B)=\mu_1(A)\mu_2(B).
\]

Для вероятностных пространств это условие всегда выполнено.

\subsubsection{Теорема 2.
Тонелли}\label{ux442ux435ux43eux440ux435ux43cux430-2.-ux442ux43eux43dux435ux43bux43bux438}

Пусть \(f\ge0\) и \(f\) измерима на \(\Omega_1\times\Omega_2\). Тогда
функции

\[
\omega_1\mapsto\int_{\Omega_2}f(\omega_1,\omega_2)\,\mu_2(d\omega_2),
\]

\[
\omega_2\mapsto\int_{\Omega_1}f(\omega_1,\omega_2)\,\mu_1(d\omega_1)
\]

измеримы, и

\[
\int f\,d(\mu_1\otimes\mu_2)
=
\int_{\Omega_1}\int_{\Omega_2}f(\omega_1,\omega_2)\,\mu_2(d\omega_2)\mu_1(d\omega_1)
\]

\[
=
\int_{\Omega_2}\int_{\Omega_1}f(\omega_1,\omega_2)\,\mu_1(d\omega_1)\mu_2(d\omega_2).
\]

Значения могут быть равны \(+\infty\).

\subsubsection{Теорема 3.
Фубини}\label{ux442ux435ux43eux440ux435ux43cux430-3.-ux444ux443ux431ux438ux43dux438}

Если \(f\in L^1(\mu_1\otimes\mu_2)\), то повторные интегралы существуют
почти всюду, являются интегрируемыми функциями и

\[
\int f\,d(\mu_1\otimes\mu_2)
=
\int_{\Omega_1}\int_{\Omega_2}f\,d\mu_2\,d\mu_1
=
\int_{\Omega_2}\int_{\Omega_1}f\,d\mu_1\,d\mu_2.
\]

\subsubsection{Следствие. Проверка условия
Фубини}\label{ux441ux43bux435ux434ux441ux442ux432ux438ux435.-ux43fux440ux43eux432ux435ux440ux43aux430-ux443ux441ux43bux43eux432ux438ux44f-ux444ux443ux431ux438ux43dux438}

Если

\[
\int_{\Omega_1}\int_{\Omega_2}|f(\omega_1,\omega_2)|\,\mu_2(d\omega_2)\mu_1(d\omega_1)<\infty,
\]

то \(f\in L^1(\mu_1\otimes\mu_2)\) и порядок интегрирования можно
менять.

\subsubsection{Теорема 4. Независимость и произведение
распределений}\label{ux442ux435ux43eux440ux435ux43cux430-4.-ux43dux435ux437ux430ux432ux438ux441ux438ux43cux43eux441ux442ux44c-ux438-ux43fux440ux43eux438ux437ux432ux435ux434ux435ux43dux438ux435-ux440ux430ux441ux43fux440ux435ux434ux435ux43bux435ux43dux438ux439}

Случайные элементы \(X\) и \(Y\) независимы тогда и только тогда, когда

\[
P_{(X,Y)}=P_X\otimes P_Y.
\]

Для \(n\) случайных элементов:

\[
P_{(X_1,\ldots,X_n)}
=
P_{X_1}\otimes\cdots\otimes P_{X_n}.
\]

\subsubsection{Теорема 5. Факторизация функций
распределения}\label{ux442ux435ux43eux440ux435ux43cux430-5.-ux444ux430ux43aux442ux43eux440ux438ux437ux430ux446ux438ux44f-ux444ux443ux43dux43aux446ux438ux439-ux440ux430ux441ux43fux440ux435ux434ux435ux43bux435ux43dux438ux44f}

Для вещественных случайных величин \(X_1,\ldots,X_n\) независимость
эквивалентна равенству

\[
F_{X_1,\ldots,X_n}(x_1,\ldots,x_n)
=
\prod_{k=1}^n F_{X_k}(x_k)
\]

для всех \(x_1,\ldots,x_n\).

\subsubsection{Теорема 6. Факторизация
плотности}\label{ux442ux435ux43eux440ux435ux43cux430-6.-ux444ux430ux43aux442ux43eux440ux438ux437ux430ux446ux438ux44f-ux43fux43bux43eux442ux43dux43eux441ux442ux438}

Если у \((X,Y)\) есть совместная плотность \(p_{X,Y}\), то \(X\) и \(Y\)
независимы тогда и только тогда, когда

\[
p_{X,Y}(x,y)=p_X(x)p_Y(y)
\]

почти всюду.

\subsubsection{Теорема 7. Ожидание произведения независимых
функций}\label{ux442ux435ux43eux440ux435ux43cux430-7.-ux43eux436ux438ux434ux430ux43dux438ux435-ux43fux440ux43eux438ux437ux432ux435ux434ux435ux43dux438ux44f-ux43dux435ux437ux430ux432ux438ux441ux438ux43cux44bux445-ux444ux443ux43dux43aux446ux438ux439}

Если \(X\) и \(Y\) независимы, то для неотрицательных измеримых \(g,h\):

\[
\mathbb E[g(X)h(Y)]
=
\mathbb E g(X)\,\mathbb E h(Y),
\]

где равенство понимается в расширенном смысле \([0,+\infty]\) через
теорему Тонелли. Случай \(0\cdot\infty\) понимается естественно: если,
например, \(\mathbb E g(X)=0\), то \(g(X)=0\) почти наверное, поэтому
обе части равны нулю. Если \(g,h\) знакопеременные, достаточно условий
\(\mathbb E|g(X)|<\infty\) и \(\mathbb E|h(Y)|<\infty\); при
независимости тогда \(g(X)h(Y)\) интегрируема и формула верна без
расширенных значений.

В частности:

\[
\mathbb E[XY]=\mathbb EX\,\mathbb EY
\]

для независимых \(X,Y\in L^1\).

\subsubsection{Теорема 8. Независимые
преобразования}\label{ux442ux435ux43eux440ux435ux43cux430-8.-ux43dux435ux437ux430ux432ux438ux441ux438ux43cux44bux435-ux43fux440ux435ux43eux431ux440ux430ux437ux43eux432ux430ux43dux438ux44f}

Если \(X\) и \(Y\) независимы, а \(g,h\) измеримы, то \(g(X)\) и
\(h(Y)\) независимы. Более общо, измеримые функции от независимых групп
случайных величин независимы.

\subsubsection{Свойство 9. Сумма независимых
величин}\label{ux441ux432ux43eux439ux441ux442ux432ux43e-9.-ux441ux443ux43cux43cux430-ux43dux435ux437ux430ux432ux438ux441ux438ux43cux44bux445-ux432ux435ux43bux438ux447ux438ux43d}

Если \(X\) и \(Y\) независимы, то

\[
P_{X+Y}=P_X*P_Y.
\]

Если есть плотности:

\[
p_{X+Y}(z)=\int_{\mathbb R}p_X(x)p_Y(z-x)\,dx.
\]

Если есть характеристические функции:

\[
\varphi_{X+Y}(t)=\varphi_X(t)\varphi_Y(t).
\]

\subsection{5. Примеры}\label{ux43fux440ux438ux43cux435ux440ux44b}

\subsubsection{Пример 1. Повторный интеграл по
прямоугольнику}\label{ux43fux440ux438ux43cux435ux440-1.-ux43fux43eux432ux442ux43eux440ux43dux44bux439-ux438ux43dux442ux435ux433ux440ux430ux43b-ux43fux43e-ux43fux440ux44fux43cux43eux443ux433ux43eux43bux44cux43dux438ux43aux443}

Пусть \(f(x,y)=x+y\) на \([0,1]\times[0,2]\) с мерой Лебега. Тогда

\[
\int_0^1\int_0^2(x+y)\,dy\,dx
=
\int_0^1(2x+2)\,dx
=
3.
\]

В другом порядке:

\[
\int_0^2\int_0^1(x+y)\,dx\,dy
=
\int_0^2\left(\frac12+y\right)\,dy
=
3.
\]

Равенство обеспечивается Фубини, так как функция ограничена на множестве
конечной меры.

\subsubsection{Пример 2. Независимые бернуллиевские
величины}\label{ux43fux440ux438ux43cux435ux440-2.-ux43dux435ux437ux430ux432ux438ux441ux438ux43cux44bux435-ux431ux435ux440ux43dux443ux43bux43bux438ux435ux432ux441ux43aux438ux435-ux432ux435ux43bux438ux447ux438ux43dux44b}

Пусть \(X,Y\) независимы и

\[
P\{X=1\}=p,\qquad P\{Y=1\}=q.
\]

Тогда

\[
P\{X=1,Y=1\}=pq,
\]

а совместное распределение задается произведениями вероятностей:

\[
P\{X=x,Y=y\}=P\{X=x\}P\{Y=y\}.
\]

Для суммы \(S=X+Y\):

\[
P\{S=2\}=pq,
\]

\[
P\{S=1\}=p(1-q)+(1-p)q,
\]

\[
P\{S=0\}=(1-p)(1-q).
\]

Это дискретная свертка.

\subsubsection{Пример 3. Независимые величины с
плотностями}\label{ux43fux440ux438ux43cux435ux440-3.-ux43dux435ux437ux430ux432ux438ux441ux438ux43cux44bux435-ux432ux435ux43bux438ux447ux438ux43dux44b-ux441-ux43fux43bux43eux442ux43dux43eux441ux442ux44fux43cux438}

Пусть \(X\) и \(Y\) независимы, причем \(X\) имеет плотность \(p_X\), а
\(Y\) - плотность \(p_Y\). Тогда

\[
p_{X,Y}(x,y)=p_X(x)p_Y(y).
\]

Например, если \(X,Y\) независимы и равномерны на \([0,1]\), то

\[
p_{X,Y}(x,y)=\mathbf 1_{[0,1]}(x)\mathbf 1_{[0,1]}(y),
\]

то есть совместная плотность постоянна на единичном квадрате.

\subsubsection{Пример 4. Плотность суммы независимых равномерных
величин}\label{ux43fux440ux438ux43cux435ux440-4.-ux43fux43bux43eux442ux43dux43eux441ux442ux44c-ux441ux443ux43cux43cux44b-ux43dux435ux437ux430ux432ux438ux441ux438ux43cux44bux445-ux440ux430ux432ux43dux43eux43cux435ux440ux43dux44bux445-ux432ux435ux43bux438ux447ux438ux43d}

Пусть \(X,Y\) независимы и равномерны на \([0,1]\). Тогда

\[
p_{X+Y}(z)
=
\int_{\mathbb R}\mathbf 1_{[0,1]}(x)\mathbf 1_{[0,1]}(z-x)\,dx.
\]

Отсюда

\[
p_{X+Y}(z)=
\begin{cases}
z, & 0\le z\le1,\\
2-z, & 1<z\le2,\\
0, & \text{иначе}.
\end{cases}
\]

Это стандартный пример свертки.

\subsubsection{Пример 5. Нулевая ковариация без
независимости}\label{ux43fux440ux438ux43cux435ux440-5.-ux43dux443ux43bux435ux432ux430ux44f-ux43aux43eux432ux430ux440ux438ux430ux446ux438ux44f-ux431ux435ux437-ux43dux435ux437ux430ux432ux438ux441ux438ux43cux43eux441ux442ux438}

Пусть \(X\) принимает значения \(-1,0,1\) с вероятностями \(1/3\), а

\[
Y=X^2.
\]

Тогда \(Y\) полностью определяется по \(X\), значит \(X\) и \(Y\)
зависимы. Но

\[
\mathbb EX=0,\qquad
\mathbb E[XY]=\mathbb E[X^3]=0,
\]

поэтому

\[
\operatorname{Cov}(X,Y)=0.
\]

Следовательно, некоррелированность не равна независимости.

\subsubsection{Пример 6. Попарная независимость не дает независимости в
совокупности}\label{ux43fux440ux438ux43cux435ux440-6.-ux43fux43eux43fux430ux440ux43dux430ux44f-ux43dux435ux437ux430ux432ux438ux441ux438ux43cux43eux441ux442ux44c-ux43dux435-ux434ux430ux435ux442-ux43dux435ux437ux430ux432ux438ux441ux438ux43cux43eux441ux442ux438-ux432-ux441ux43eux432ux43eux43aux443ux43fux43dux43eux441ux442ux438}

Пусть бросаются две независимые честные монеты, а события:

\[
A=\{\text{первая монета - орел}\},
\]

\[
B=\{\text{вторая монета - орел}\},
\]

\[
C=\{\text{результаты монет совпали}\}.
\]

Тогда \(A,B,C\) попарно независимы, но

\[
P(A\cap B\cap C)=\frac14
\]

и

\[
P(A)P(B)P(C)=\frac18.
\]

Значит независимости в совокупности нет.

\subsection{6. Схемы доказательств /
обоснований}\label{ux441ux445ux435ux43cux44b-ux434ux43eux43aux430ux437ux430ux442ux435ux43bux44cux441ux442ux432-ux43eux431ux43eux441ux43dux43eux432ux430ux43dux438ux439}

\textbf{Схема доказательства существования произведения мер.}

Сначала на прямоугольниках задают

\[
\mu(A\times B)=\mu_1(A)\mu_2(B).
\]

Затем продолжают эту функцию с полукольца или алгебры прямоугольников на
порожденную \(\sigma\)-алгебру. Единственность получается из теоремы
продолжения меры при \(\sigma\)-конечности. Для вероятностных мер это
особенно просто, потому что меры конечны.

\textbf{Схема доказательства Тонелли.}

\begin{enumerate}
\def\labelenumi{\arabic{enumi}.}
\tightlist
\item
  Сначала проверяют формулу для индикатора прямоугольника:
\end{enumerate}

\[
f=\mathbf 1_{A\times B}.
\]

Тогда обе стороны равны \(\mu_1(A)\mu_2(B)\).

\begin{enumerate}
\def\labelenumi{\arabic{enumi}.}
\setcounter{enumi}{1}
\item
  Затем переходят к индикаторам множеств из
  \(\mathcal F_1\otimes\mathcal F_2\) через монотонный класс.
\item
  Потом доказывают формулу для неотрицательных простых функций по
  линейности.
\item
  Для произвольной \(f\ge0\) берут простые функции \(f_n\uparrow f\) и
  применяют теорему Беппо Леви о монотонной сходимости из Билет 07.
\end{enumerate}

\textbf{Схема доказательства Фубини.}

Для интегрируемой функции \(f\) раскладывают

\[
f=f^+-f^-.
\]

Так как

\[
\int |f|\,d(\mu_1\otimes\mu_2)<\infty,
\]

то интегралы от \(f^+\) и \(f^-\) конечны. К ним применяют Тонелли,
затем вычитают равенства. Это дает существование повторных интегралов
почти всюду и равенство всех трех интегралов.

\textbf{Схема доказательства критерия независимости через произведение
распределений.}

Если \(X\) и \(Y\) независимы, то для прямоугольников \(A\times B\):

\[
P_{(X,Y)}(A\times B)
=
P\{X\in A,Y\in B\}
=
P_X(A)P_Y(B)
=
(P_X\otimes P_Y)(A\times B).
\]

Так как прямоугольники порождают произведенную \(\sigma\)-алгебру,
равенство мер на прямоугольниках переносится на всю \(\sigma\)-алгебру.
Обратное следует сразу из той же формулы на прямоугольниках.

\textbf{Схема доказательства формулы
\(\mathbb E[g(X)h(Y)]=\mathbb E g(X)\mathbb E h(Y)\).}

По теореме о замене переменной:

\[
\mathbb E[g(X)h(Y)]
=
\int g(x)h(y)\,P_{(X,Y)}(dx,dy).
\]

По независимости:

\[
P_{(X,Y)}=P_X\otimes P_Y.
\]

По Фубини:

\[
\int g(x)h(y)\,P_X(dx)P_Y(dy)
=
\left(\int g\,dP_X\right)\left(\int h\,dP_Y\right).
\]

\subsection{7. Что написать на
доске}\label{ux447ux442ux43e-ux43dux430ux43fux438ux441ux430ux442ux44c-ux43dux430-ux434ux43eux441ux43aux435}

\begin{enumerate}
\def\labelenumi{\arabic{enumi}.}
\tightlist
\item
  Произведенная мера:
\end{enumerate}

\[
(\mu_1\otimes\mu_2)(A\times B)=\mu_1(A)\mu_2(B).
\]

\begin{enumerate}
\def\labelenumi{\arabic{enumi}.}
\setcounter{enumi}{1}
\tightlist
\item
  Тонелли:
\end{enumerate}

\[
f\ge0
\quad\Rightarrow\quad
\int f\,d(\mu_1\otimes\mu_2)
=
\int\int f\,d\mu_2\,d\mu_1
=
\int\int f\,d\mu_1\,d\mu_2.
\]

\begin{enumerate}
\def\labelenumi{\arabic{enumi}.}
\setcounter{enumi}{2}
\tightlist
\item
  Фубини:
\end{enumerate}

\[
f\in L^1(\mu_1\otimes\mu_2)
\quad\Rightarrow\quad
\int f\,d(\mu_1\otimes\mu_2)
=
\int\int f\,d\mu_2\,d\mu_1.
\]

\begin{enumerate}
\def\labelenumi{\arabic{enumi}.}
\setcounter{enumi}{3}
\tightlist
\item
  Независимость:
\end{enumerate}

\[
P(A\cap B)=P(A)P(B).
\]

\begin{enumerate}
\def\labelenumi{\arabic{enumi}.}
\setcounter{enumi}{4}
\tightlist
\item
  Независимость случайных элементов:
\end{enumerate}

\[
P_{(X,Y)}=P_X\otimes P_Y.
\]

\begin{enumerate}
\def\labelenumi{\arabic{enumi}.}
\setcounter{enumi}{5}
\tightlist
\item
  Ожидание произведения:
\end{enumerate}

\[
\mathbb E[g(X)h(Y)]
=
\mathbb E g(X)\,\mathbb E h(Y).
\]

\begin{enumerate}
\def\labelenumi{\arabic{enumi}.}
\setcounter{enumi}{6}
\tightlist
\item
  Плотности и свертка:
\end{enumerate}

\[
p_{X,Y}(x,y)=p_X(x)p_Y(y),
\qquad
p_{X+Y}(z)=\int p_X(x)p_Y(z-x)\,dx.
\]

\subsection{8. Типичные дополнительные
вопросы}\label{ux442ux438ux43fux438ux447ux43dux44bux435-ux434ux43eux43fux43eux43bux43dux438ux442ux435ux43bux44cux43dux44bux435-ux432ux43eux43fux440ux43eux441ux44b}

\textbf{Вопрос. Чем Тонелли отличается от Фубини?}

Тонелли работает для \(f\ge0\) и допускает бесконечные значения
интегралов. Фубини работает для знакопеременных функций, но требует
\(f\in L^1\), то есть \(\int |f|<\infty\).

\textbf{Вопрос. Можно ли менять порядок интегрирования без условий?}

Нет. Для неотрицательных функций можно по Тонелли. Для знакопеременных
нужно проверять абсолютную интегрируемость или другое условие,
гарантирующее применимость Фубини.

\textbf{Вопрос. Почему в Фубини повторные интегралы определены только
почти всюду?}

Потому что для отдельных сечений интеграл может быть не определен или
бесконечен на множестве внешней переменной меры ноль. Интеграл Лебега не
различает изменения на нулевых множествах.

\textbf{Вопрос. Что значит независимость случайных величин через
\(\sigma\)-алгебры?}

Это независимость всей информации, содержащейся в величинах. Формально
независимы \(\sigma(X)\) и \(\sigma(Y)\), то есть все события вида
\(\{X\in A\}\) и \(\{Y\in B\}\).

\textbf{Вопрос. Как связаны независимость и совместное распределение?}

Независимость означает, что совместный закон равен произведению
маргинальных законов:

\[
P_{(X,Y)}=P_X\otimes P_Y.
\]

\textbf{Вопрос. Как проверить независимость по функции распределения?}

Для вещественных \(X,Y\) нужно проверить

\[
F_{X,Y}(x,y)=F_X(x)F_Y(y)
\]

для всех \(x,y\).

\textbf{Вопрос. Как проверить независимость по плотности?}

Если есть совместная плотность, то нужно проверить

\[
p_{X,Y}(x,y)=p_X(x)p_Y(y)
\]

почти всюду.

\textbf{Вопрос. Почему из независимости следует
\(\mathbb E[XY]=\mathbb EX\,\mathbb EY\)?}

Потому что совместное распределение равно \(P_X\otimes P_Y\), а по
Фубини интеграл от \(xy\) по произведению мер раскладывается в
произведение интегралов.

\textbf{Вопрос. Верно ли обратное: если
\(\mathbb E[XY]=\mathbb EX\,\mathbb EY\), то \(X,Y\) независимы?}

Нет. Это означает только нулевую ковариацию при наличии вторых моментов.
Независимость сильнее.

\textbf{Вопрос. Почему попарная независимость не равна независимости в
совокупности?}

Потому что попарная независимость проверяет только пересечения двух
событий, а независимость в совокупности требует факторизации всех
пересечений любых поднаборов.

\textbf{Вопрос. Что происходит с функциями от независимых величин?}

Если \(X\) и \(Y\) независимы, а \(g,h\) измеримы, то \(g(X)\) и
\(h(Y)\) тоже независимы.

\textbf{Вопрос. Как распределение суммы связано с независимостью?}

При независимости закон суммы есть свертка:

\[
P_{X+Y}=P_X*P_Y.
\]

Без независимости маргинальных распределений недостаточно: нужен
совместный закон.

\subsection{9. Типичные
ошибки}\label{ux442ux438ux43fux438ux447ux43dux44bux435-ux43eux448ux438ux431ux43aux438}

\begin{enumerate}
\def\labelenumi{\arabic{enumi}.}
\tightlist
\item
  Менять порядок интегрирования для знакопеременной функции без проверки
  \(\int |f|<\infty\).
\item
  Называть любую перестановку интегралов ``Фубини'', хотя для \(f\ge0\)
  точнее говорить ``Тонелли''.
\item
  Забывать, что повторные интегралы в Фубини определены почти всюду, а
  не обязательно при каждом значении внешней переменной.
\item
  Путать попарную независимость и независимость в совокупности.
\item
  Считать, что некоррелированность означает независимость.
\item
  Проверять независимость только на одном значении функции
  распределения, а не для всех \(x,y\).
\item
  Забывать ``почти всюду'' в равенстве плотностей
\end{enumerate}

\[
p_{X,Y}=p_Xp_Y.
\]

\begin{enumerate}
\def\labelenumi{\arabic{enumi}.}
\setcounter{enumi}{7}
\tightlist
\item
  Использовать формулу
\end{enumerate}

\[
\mathbb E[g(X)h(Y)]=\mathbb E g(X)\mathbb E h(Y)
\]

без независимости или без условий существования ожиданий. 9. Считать,
что распределение суммы всегда является сверткой маргинальных
распределений. Это верно только при независимости. 10. Путать
произведение случайных величин \(XY\) и произведение мер
\(P_X\otimes P_Y\).

\subsection{10. Связи с другими
билетами}\label{ux441ux432ux44fux437ux438-ux441-ux434ux440ux443ux433ux438ux43cux438-ux431ux438ux43bux435ux442ux430ux43cux438}

\textbf{Билет 06.} Совместные распределения, конечномерные распределения
и плотности дают язык, в котором формулируется независимость.

\textbf{Билет 07.} Тонелли и Фубини опираются на интеграл Лебега и
теоремы о монотонной и мажорированной сходимости.

\textbf{Билет 08.} Для независимых сумм характеристические функции
перемножаются:

\[
\varphi_{X+Y}=\varphi_X\varphi_Y.
\]

\textbf{Билет 10.} Условное ожидание использует независимость от
\(\sigma\)-алгебры: если \(X\) независима от \(\mathcal G\), то

\[
\mathbb E[X\mid\mathcal G]=\mathbb EX.
\]

\textbf{Билет 11.} Замена переменной в интеграле используется вместе с
Фубини для перехода от \(\mathbb E g(X,Y)\) к интегралу по
\(P_{(X,Y)}\).

\textbf{Билет 14.} Независимость дает нулевую ковариацию, но нулевая
ковариация не дает независимости.

\textbf{Билет 18.} Для совместно гауссовских величин некоррелированность
эквивалентна независимости.

\textbf{Билет 19.} Независимые приращения и блуждания используют
факторизацию совместных законов и свертки.

\textbf{Билет 21.} Закон больших чисел и центральная предельная теорема
требуют независимости или ее ослабленных вариантов.

\subsection{11. Минимум на
удовлетворительно}\label{ux43cux438ux43dux438ux43cux443ux43c-ux43dux430-ux443ux434ux43eux432ux43bux435ux442ux432ux43eux440ux438ux442ux435ux43bux44cux43dux43e}

Нужно уметь сказать:

\begin{enumerate}
\def\labelenumi{\arabic{enumi}.}
\tightlist
\item
  Теорема Фубини позволяет считать интеграл по произведению мер как
  повторный интеграл.
\item
  Для \(f\ge0\) работает Тонелли, для знакопеременной \(f\) нужно
  \(f\in L^1\).
\item
  Независимость событий:
\end{enumerate}

\[
P(A\cap B)=P(A)P(B).
\]

\begin{enumerate}
\def\labelenumi{\arabic{enumi}.}
\setcounter{enumi}{3}
\tightlist
\item
  Независимость случайных величин означает
\end{enumerate}

\[
P_{(X,Y)}=P_X\otimes P_Y.
\]

\begin{enumerate}
\def\labelenumi{\arabic{enumi}.}
\setcounter{enumi}{4}
\tightlist
\item
  При независимости
\end{enumerate}

\[
\mathbb E[XY]=\mathbb EX\,\mathbb EY.
\]

\subsection{12. Минимум на
хорошо}\label{ux43cux438ux43dux438ux43cux443ux43c-ux43dux430-ux445ux43eux440ux43eux448ux43e}

Помимо минимума, нужно:

\begin{enumerate}
\def\labelenumi{\arabic{enumi}.}
\tightlist
\item
  различать Тонелли и Фубини;
\item
  записать условие абсолютной интегрируемости;
\item
  сформулировать независимость \(\sigma\)-алгебр и случайных величин;
\item
  объяснить факторизацию функции распределения:
\end{enumerate}

\[
F_{X,Y}(x,y)=F_X(x)F_Y(y);
\]

\begin{enumerate}
\def\labelenumi{\arabic{enumi}.}
\setcounter{enumi}{4}
\tightlist
\item
  привести пример свертки для суммы независимых величин;
\item
  сказать, что попарная независимость не равна независимости в
  совокупности.
\end{enumerate}

\subsection{13. Что добавить для
отлично}\label{ux447ux442ux43e-ux434ux43eux431ux430ux432ux438ux442ux44c-ux434ux43bux44f-ux43eux442ux43bux438ux447ux43dux43e}

Для сильного ответа стоит добавить:

\begin{enumerate}
\def\labelenumi{\arabic{enumi}.}
\tightlist
\item
  произведенную \(\sigma\)-алгебру \(\mathcal F_1\otimes\mathcal F_2\);
\item
  существование и единственность произведенной меры при
  \(\sigma\)-конечности;
\item
  доказательную схему через индикаторы прямоугольников, простые функции
  и монотонную сходимость;
\item
  равенство плотностей почти всюду:
\end{enumerate}

\[
p_{X,Y}=p_Xp_Y;
\]

\begin{enumerate}
\def\labelenumi{\arabic{enumi}.}
\setcounter{enumi}{4}
\tightlist
\item
  формулу свертки плотностей;
\item
  контрпример ``некоррелированность не означает независимость'';
\item
  контрпример ``попарная независимость не означает независимость в
  совокупности''.
\end{enumerate}

\subsection{14. Устная версия ответа на 5
минут}\label{ux443ux441ux442ux43dux430ux44f-ux432ux435ux440ux441ux438ux44f-ux43eux442ux432ux435ux442ux430-ux43dux430-5-ux43cux438ux43dux443ux442}

Сначала я ввожу произведение измеримых пространств. Если есть
\((\Omega_1,\mathcal F_1,\mu_1)\) и \((\Omega_2,\mathcal F_2,\mu_2)\),
то на произведении берут

\[
\mathcal F_1\otimes\mathcal F_2
=
\sigma\{A\times B\}.
\]

При \(\sigma\)-конечности существует произведенная мера:

\[
(\mu_1\otimes\mu_2)(A\times B)=\mu_1(A)\mu_2(B).
\]

Дальше формулирую Тонелли: если \(f\ge0\) измерима, то

\[
\int f\,d(\mu_1\otimes\mu_2)
=
\int\int f\,d\mu_2\,d\mu_1
=
\int\int f\,d\mu_1\,d\mu_2.
\]

Значение может быть бесконечным. Затем формулирую Фубини: если

\[
\int |f|\,d(\mu_1\otimes\mu_2)<\infty,
\]

то повторные интегралы существуют почти всюду, конечны после внешнего
интегрирования и равны интегралу по произведению. Отличие такое: Тонелли
- для неотрицательных функций без проверки конечности, Фубини - для
интегрируемых функций.

После этого перехожу к независимости. События \(A\) и \(B\) независимы,
если

\[
P(A\cap B)=P(A)P(B).
\]

\(\sigma\)-алгебры независимы, если такая факторизация верна для всех
событий из них. Случайные величины независимы, если независимы
порожденные ими \(\sigma\)-алгебры. Эквивалентная и самая полезная
запись:

\[
P_{(X,Y)}=P_X\otimes P_Y.
\]

Для вещественных величин это значит

\[
F_{X,Y}(x,y)=F_X(x)F_Y(y).
\]

Если есть плотности, то

\[
p_{X,Y}(x,y)=p_X(x)p_Y(y)
\]

почти всюду.

Главное вычислительное следствие: если \(X,Y\) независимы, то

\[
\mathbb E[g(X)h(Y)]
=
\mathbb E g(X)\mathbb E h(Y),
\]

потому что интеграл идет по произведенной мере и раскладывается по
Фубини. В частности, для независимых интегрируемых \(X,Y\):

\[
\mathbb E[XY]=\mathbb EX\,\mathbb EY.
\]

Для суммы независимых величин получаем свертку:

\[
p_{X+Y}(z)=\int p_X(x)p_Y(z-x)\,dx,
\]

а для характеристических функций:

\[
\varphi_{X+Y}(t)=\varphi_X(t)\varphi_Y(t).
\]

В конце называю ловушки: нельзя менять порядок интегрирования без
условий, попарная независимость не равна независимости в совокупности, а
нулевая ковариация не означает независимости.

\subsection{15. Если осталось 2 минуты перед
экзаменом}\label{ux435ux441ux43bux438-ux43eux441ux442ux430ux43bux43eux441ux44c-2-ux43cux438ux43dux443ux442ux44b-ux43fux435ux440ux435ux434-ux44dux43aux437ux430ux43cux435ux43dux43eux43c}

Запомнить три блока.

\textbf{Фубини/Тонелли:}

\[
f\ge0
\Rightarrow
\int f\,d(\mu\otimes\nu)
=
\int\int f(x,y)\,\nu(dy)\mu(dx)
=
\int\int f(x,y)\,\mu(dx)\nu(dy).
\]

Для знакопеременной функции нужно

\[
\int |f|\,d(\mu\otimes\nu)<\infty.
\]

\textbf{Независимость:}

\[
P(A\cap B)=P(A)P(B),
\qquad
P_{(X,Y)}=P_X\otimes P_Y.
\]

\textbf{Главные следствия:}

\[
\mathbb E[g(X)h(Y)]
=
\mathbb E g(X)\mathbb E h(Y),
\]

\[
p_{X,Y}=p_Xp_Y,
\qquad
P_{X+Y}=P_X*P_Y.
\]

Главная ловушка: Фубини требует условий, а независимость намного сильнее
некоррелированности.

\newpage

\end{document}
